% Generated draft: requires mathematical and source review.
\documentclass[11pt]{article}
\usepackage[margin=25mm]{geometry}
\usepackage{fontspec}
\setmainfont{Latin Modern Roman}
\usepackage{amsmath,amssymb,mathtools,unicode-math}
\setmathfont{Latin Modern Math}
\usepackage{enumitem,needspace,xurl,hyperref}
\setlength{\parindent}{0pt}
\setlength{\parskip}{5pt}
\setlength{\emergencystretch}{3em}
\newcommand{\N}{\mathbb{N}}
\newcommand{\Z}{\mathbb{Z}}
\newcommand{\Q}{\mathbb{Q}}
\newcommand{\R}{\mathbb{R}}
\newcommand{\C}{\mathbb{C}}
\newcommand{\F}{\mathbb{F}}
\newcommand{\A}{\mathbb{A}}
\newcommand{\PP}{\mathbb{P}}
\newcommand{\E}{\mathbb{E}}
\newcommand{\eps}{\varepsilon}
\newcommand{\dd}{\,\mathrm{d}}
\begin{document}
\section*{UnsolvedMath ID: 800035}
\textbf{Source problem code:} catalogue entry
\subsection{Definitions and mathematical statement}
A semigroup is a nonempty set with an associative multiplication. A group \(G\) is periodic if every element has finite order, that is,
\[
\forall g\in G\ \exists n\in\Z_{\geq 1}\quad g^n=1_G.
\]
No common bound on these orders is assumed. For subsets \(A,B\subseteq G\), define their setwise product by
\[
AB=\{ab:a\in A,\ b\in B\}.
\]
Both \(\mathcal P(G)=\{A:A\subseteq G\}\) and \(\mathcal P_+(G)=\mathcal P(G)\setminus\{\varnothing\}\) are semigroups under this operation. The supplied record does not specify whether the empty subset belongs to the intended power semigroup. Write \(\mathscr P(G)\) for the intended choice, using the same convention throughout the statement below; no equivalence between the two resulting formulations is asserted here.

For a finite nonempty subset \(F\) of a semigroup \(S\), let
\[
\langle F\rangle=\{f_1\cdots f_k:k\in\Z_{\geq1},\ f_1,\ldots,f_k\in F\}
\]
be the subsemigroup it generates. An embedding into \(\mathscr P(G)\) means an injective map \(\phi\) satisfying \(\phi(xy)=\phi(x)\phi(y)\); no preservation of an identity is required.

\textbf{Question (S. G. Bershadsky and S. I. Kublanovsky; Problem 37).} Is the following implication true for every semigroup \(S\)? If, for every finite nonempty \(F\subseteq S\), there exist a periodic group \(G_F\) and an injective semigroup homomorphism
\[
\phi_F:\langle F\rangle\longrightarrow\mathscr P(G_F),
\]
then there exist a periodic group \(G\) and an injective semigroup homomorphism
\[
\phi:S\longrightarrow\mathscr P(G).
\]
The groups and embeddings in the hypothesis may depend on \(F\); no compatibility between these embeddings is assumed.
\subsection{Short English statement}
If every finitely generated subsemigroup of a semigroup embeds into the power semigroup of some periodic group, must the entire semigroup embed into the power semigroup of a periodic group? The convention concerning the empty subset remains unspecified in the supplied source.
\subsection{Sources}
\begin{itemize}
\item Supplied local UnsolvedMath registry record: permanent ID \(800035\), exact source problem code \texttt{catalogue entry}, title \emph{The Lyapin's Notebook --- Problem 37}; upstream revision \texttt{d371ee917ac9647fbb06cd6fc529deec6f3dda82}.
\item Supplied source attribution: Bershadsky, Kublanovsky, and Mashevitzky, \emph{The Lyapin's Notebook: Unsolved Problems in Semigroup Theory} (2026), Problem 37, \url{https://arxiv.org/abs/2604.04763}. The problem statement specifically attributes the question to S. G. Bershadsky and S. I. Kublanovsky; the registry's broader \texttt{proposed\_by} field lists Bershadsky, Kublanovsky, and Mashevitzky.
\item The supplied provenance records access on 2026-10-16 and extraction by \texttt{arxiv\_latex}. These are supplied provenance claims, not independently verified access or extraction. The registry labels the problem \texttt{open}; that label is supplied curation, not an independently verified literature assessment.
\end{itemize}
\end{document}
